\section{The descendent broken-line correspondence}
\label{sec:descendent-correspondence}

We now construct the class in
Theorem~\ref{thm:intro-descendent} and the cellular model of
Definition~\ref{def:intro-relative-pl-bordism}.  Fix a finite labelled window
$S=(I,P,\mathcal E)$ as in
Definition~\ref{def:intro-finite-window}, a deformation order $N$, the finite
root stage, and a framed compact annular corridor defined there.  All coefficients are
taken in the finite Artin ring $R_{S,N}$ of
\eqref{eq:intro-window-ring}.  The closure conditions on $\mathcal E$ retain
every proper support, every open boundary graph and bundle map in
\eqref{eq:intro-exact-open-datum}, and every closed-component splitting,
including those whose Euler rank does not equal the dimension of its compact
moduli factor.
All spaces below are therefore finite-dimensional and compact.  The virtual
chains are rational PL branched zero loci of finite orbibundle sections; no
analytic infinite-window virtual class is used.

The construction has four parts.  First, a relative twisted de~Rham complex
for the Fermat polynomial supplies a residue cocycle compatible with the two
coordinate divisors.  Second, a deformation retract of this complex gives
coherent primitives for the factorization faces of tropical forests.  Third,
the Witten-bundle splitting realizes insertion of a cotangent class by a
parameterized Euler chain.  Finally, these data are assembled on a compactified
incidence space of tropical forests and $W$-spin curves.  The boundary
stratification of that compactification gives the two cotangent recursions by
Stokes' formula.

The chain-level argument has the following conceptual order:
\[
\begin{tikzcd}[column sep=2.0em,row sep=1.7em]
\text{relative residue complex}\arrow[r] &
\text{forest primitives}\arrow[r] &
\text{incidence zero chain}\arrow[d] \\
\text{GKT moments} &
\text{cotangent recursions}\arrow[l] &
\text{boundary strata}\arrow[l]
\end{tikzcd}
\]
Here a broken-line chain is a cellular chain whose strata parametrize
decorated tropical forests together with their $W$-spin data; it is not a
single Gross--Siebert broken line.

\subsection{Relative contact Stokes}

Put $W_0=X^4+Y^4$, $\Omega=dX\wedge dY$, and
$D_0=\hbar d+dW_0\wedge$.  Let
$\mathsf D_x=\{X=0\}$, $\mathsf D_y=\{Y=0\}$, and
$\mathsf D_{xy}=\mathsf D_x\cap\mathsf D_y$.  On the two divisors and their
intersection put
\[
 D_x=\hbar d+d(Y^4)\wedge,\qquad
 D_y=\hbar d+d(X^4)\wedge,\qquad D_{xy}=0.
\]
Let $\mathcal B_A,\mathcal B_x,\mathcal B_y,$ and
$\mathcal B_{xy}=\QQ[\hbar]$ be the corresponding twisted polynomial
de~Rham complexes, and let $r_x,r_y,r_{xy}^x,r_{xy}^y$ be their ordinary
restriction maps.  Define
\begin{equation}
 \mathcal B_{\mathrm{con}}^k
 =\mathcal B_A^k\oplus\mathcal B_x^{k-1}
  \oplus\mathcal B_y^{k-1}\oplus\mathcal B_{xy}^{k-2}
 \label{eq:des-contact-complex}
\end{equation}
with differential
\begin{equation}
\begin{aligned}
 \delta(\omega,\eta_x,\eta_y,\xi)
 =\bigl(&D_0\omega,r_x\omega-D_x\eta_x,r_y\omega-D_y\eta_y,\\
 &r_{xy}^x\eta_x-r_{xy}^y\eta_y+D_{xy}\xi\bigr).
\end{aligned}
\label{eq:des-contact-differential}
\end{equation}
The restriction square commutes, so $\delta^2=0$.

For $0\leq r\leq3$, set
\begin{equation}
 G_r(p)=\prod_{j=0}^{p-1}\frac{r+1+4j}{4},\qquad G_r(0)=1.
 \label{eq:des-gamma-weight}
\end{equation}
The bulk, divisor, and corner traces are
\begin{equation}
\begin{aligned}
 \tau_A(X^{r+4p}Y^{s+4q}\Omega)&=(-\hbar)^{p+q}G_r(p)G_s(q),\\
 \tau_x(Y^{s+4q}dY)&=4(-\hbar)^qG_s(q),\\
 \tau_y(X^{r+4p}dX)&=-4(-\hbar)^pG_r(p),\qquad
 \tau_{xy}(c)=16c.
\end{aligned}
\label{eq:des-contact-traces}
\end{equation}
The factors $4$ and $16$ are the inverse traces of $B\mu _4$ and
$B(\mu _4^2)$.  Set
\begin{equation}
 \operatorname{Tr}_{\mathrm{con}}(\omega,\eta_x,\eta_y,\xi)
 =\tau_A(\omega)-\tau_x(\eta_x)-\tau_y(\eta_y)-\tau_{xy}(\xi).
 \label{eq:des-total-trace}
\end{equation}

Let $\widetilde R_{S,N}$ be the finite $R_{S,N}$-algebra obtained by
adjoining the roots occurring at the chosen root stage.  Fix a chamber
$\Gamma_0$ in which the developed coordinates are $(X,Y)$ and the reference
polynomial is $W_0=X^4+Y^4$; this is the \emph{reference Fermat chamber}.
For any chamber frame $\Gamma$, let
$\mathsf T_{\Gamma\Gamma_0}$ be the ordered composition of the wall and
root-chart maps along its framed path, and define
\begin{equation}
 \mathcal B_{\Gamma,\mathrm{con}}
 :=\mathsf T_{\Gamma\Gamma_0}
   (\mathcal B_{\mathrm{con}}\otimes_{\QQ}
    \widetilde R_{S,N}),\quad
 \delta_\Gamma:=\mathsf T_{\Gamma\Gamma_0}\delta
  \mathsf T_{\Gamma\Gamma_0}^{-1},\quad
 \operatorname{Tr}_\Gamma:=\operatorname{Tr}_{\mathrm{con}}
  \mathsf T_{\Gamma\Gamma_0}^{-1}.
 \label{eq:des-transported-contact-complex}
\end{equation}
Thus Laurent and root-valued expressions live in the transported complex,
not in the original polynomial presentation.  If $\Gamma'$ is another
frame, then
$\mathsf T_{\Gamma'\Gamma}:=\mathsf T_{\Gamma'\Gamma_0}
\mathsf T_{\Gamma\Gamma_0}^{-1}$ is a chain isomorphism and
\begin{equation}
 \operatorname{Tr}_{\Gamma'}\mathsf T_{\Gamma'\Gamma}
 =\operatorname{Tr}_\Gamma.
 \label{eq:des-trace-transport}
\end{equation}
The transported coordinate divisors are part of
$\mathcal B_{\Gamma,\mathrm{con}}$; no wall map is required to preserve the
original pair $\{X=0\}\cup\{Y=0\}$.

\begin{proposition}[Relative contact Stokes formula]
\label{prop:des-contact-stokes}
One has $\operatorname{Tr}_{\mathrm{con}}\delta=0$.  Moreover,
\eqref{eq:des-contact-complex} has a coordinate-equivariant deformation
retract onto
\begin{equation}
 \bigoplus_{0\leq r,s\leq3}\QQ[\hbar]X^rY^s\Omega.
 \label{eq:des-relative-normal-form}
\end{equation}
\end{proposition}

\begin{proof}
The recurrence
$G_r(p+1)=\frac{r+1+4p}{4}G_r(p)$ gives
\begin{equation}
 \tau_A(D_0H)=\tau_x(r_xH)+\tau_y(r_yH),\quad
 \tau_x(D_xf)=16f(0),\quad \tau_y(D_yg)=-16g(0).
 \label{eq:des-integration-by-parts}
\end{equation}
Substitution in \eqref{eq:des-contact-differential} proves the Stokes
identity.  In one variable, let $R_4$ send $t^{r+4p}dt$ to
$(-\hbar)^pG_r(p)t^r dt$, and put
$R_{\mathrm{rel}}(gdt,c\epsilon)=R_4(gdt)-4ct^3dt$.  The recursive primitive
\begin{equation}
 s_4(t^a)=\frac14t^{a-3}
 -\hbar\frac{a-3}{4}s_4(t^{a-4})\quad(a\geq4),
 \label{eq:des-one-variable-homotopy}
\end{equation}
with $s_4(t^a)=0$ for $a\leq3$, satisfies the contraction identity.  The
full relative homotopy sends $(gdt,c\epsilon)$ to $s_4(g)+c$; together with
the displayed reduction it satisfies the degreewise contraction identity
$\delta h_{\mathrm{rel}}+h_{\mathrm{rel}}\delta
=1-\iota R_{\mathrm{rel}}$.  The oriented tensor product of the two contractions gives
\eqref{eq:des-relative-normal-form}; averaging its two orders makes it
coordinate equivariant.
\end{proof}

This relative complex is necessary.  Although
$D_0(\frac14dY)=X^3\Omega$ in the plane complex, in
\eqref{eq:des-contact-complex} the complete boundary is
\begin{equation}
 \delta(\tfrac14dY,0,0,0)
 =(X^3\Omega,\tfrac14dY,0,0).
 \label{eq:des-residue-three}
\end{equation}
The bulk and divisor traces both equal one and cancel.  Thus residue-three
contact data survive without assigning a nonzero trace to an ordinary exact
form.

\subsection{Factorization chains for tropical forests}

For a finite nonempty set $A$, put
$V_A=\ker(\QQ^A\xrightarrow{\sum}\QQ)$.  If $F$ is a rooted forest with
leaves $I$, write $R(F)$ for its root clusters and
$\operatorname{In}(v)$ for the child clusters entering an internal vertex.
The block-sum exact sequences
\begin{equation}
 0\longrightarrow\bigoplus_{B\in\operatorname{In}(v)}V_B
 \longrightarrow V_{\coprod B}\longrightarrow
 V_{\operatorname{In}(v)}\longrightarrow0
 \label{eq:des-block-sum}
\end{equation}
give a determinant isomorphism
\begin{equation}
 \mathfrak o_F=
 \det\left(V_{R(F)}\oplus\bigoplus_{v}V_{\operatorname{In}(v)}\right)
 \xrightarrow{\ \Phi_F\ }\det V_I.
 \label{eq:des-forest-orientation}
\end{equation}
Children, vertices, and roots are ordered by least label.  Cutting two edges
in opposite orders changes $\Phi_F$ by the exchange of the two outward
metric normals, exactly the cubical boundary sign.

For a nonempty cluster $C\subseteq I$, use the twist representatives fixed
in Definition~\ref{def:intro-finite-window} and write
\begin{equation}
 \sum_{i\in C}a_i=r_C+4\ell_C,\qquad
 \sum_{i\in C}b_i=s_C+4m_C,\qquad 0\leq r_C,s_C\leq3.
 \label{eq:des-cluster-residues}
\end{equation}
Its balanced exponent at position $p$ is
\begin{equation}
 \mathbf b_{C,p}=(r_C+4p,s_C+4(N_C-p)),\qquad 0\leq p\leq N_C.
 \label{eq:des-cluster-profile}
\end{equation}
Every vertex $v$ of a decorated forest records a cluster and either one of
these profiles or the exponent of a specified wall Hamiltonian.  Denote that
explicit exponent by
$\mathbf b_v=(b_{v,x},b_{v,y})\in\NN^2$.  Its scalar
$c_v\in\widetilde R_{S,N}$ is the displayed coefficient of that Hamiltonian
or of the normalized fourth-root collar.  A closed $W$-spin factor is not
put into $c_v$: it remains a geometric Euler factor in the incidence space
below.  Thus, beyond the two specified singleton Hamiltonian normalizations,
no higher-support GKT open invariant, endpoint coefficient, or
cotangent-recursion value enters a vertex scalar.

A \emph{decorated tropical forest} $h$ is a stable labelled forest together with
these choices at all of its vertices and the ordered wall word on each edge.
The coefficient associated with $h$ is the single bulk form
\begin{equation}
 \beta_h=c_hX^{A_h}Y^{B_h}\Omega,\qquad
 c_h=\prod_vc_v,\qquad
 (A_h,B_h)=\sum_v\mathbf b_v.
 \label{eq:des-single-seed}
\end{equation}
All wall, contact, and homotopy operators thereafter act linearly on this
one element of $\mathcal B_{\mathrm{con}}$; forms from distinct vertices are
never multiplied.

For $\mathbf b=(r+4p,s+4q)$ with $0\leq r,s\leq3$, put
\begin{equation}
 w(\mathbf b)=G_r(p)G_s(q).
 \label{eq:des-contact-weight}
\end{equation}
Let $\pi$ be a partition of a label set $J$, with $a=|\pi|$.  Choose one
profile \eqref{eq:des-cluster-profile} on each block and set
$\mathbf b_\pi=\sum_{C\in\pi}\mathbf b_C$.  The parent profile
$\mathbf b_J$ has the same two residues and the factorization degree one
lower at each new forest edge; hence uniquely
$\mathbf b_\pi=\mathbf b_J+4(c_x,c_y)$ with
$c_x+c_y=a-1$.  Let $c$ be the product of all vertex scalars outside this
contracted corolla and of the corolla's displayed wall coefficients.  Define
\begin{equation}
\begin{aligned}
 \lambda_{J,\pi}&=-\frac{w(\mathbf b_\pi)}{w(\mathbf b_J)},\\
 \beta^{\mathrm{sp}}_{J,\pi}
   &=cX^{(\mathbf b_\pi)_x}Y^{(\mathbf b_\pi)_y}\Omega,\\
 \beta^{\mathrm{par}}_{J,\pi}
   &=c(-\hbar)^{a-1}\lambda_{J,\pi}
     X^{(\mathbf b_J)_x}Y^{(\mathbf b_J)_y}\Omega.
\end{aligned}
 \label{eq:des-split-parent-seeds}
\end{equation}
If $\mathbf b_J=(r+4p,s+4q)$, the Gamma recurrence gives the explicit
cancellation
\begin{equation}
 \operatorname{Tr}_{\mathrm{con}}
 \bigl(\beta^{\mathrm{sp}}_{J,\pi}
      +\beta^{\mathrm{par}}_{J,\pi},0,0,0\bigr)
 =c(-\hbar)^{p+q+a-1}
 \bigl(w(\mathbf b_\pi)+\lambda_{J,\pi}w(\mathbf b_J)\bigr)=0.
 \label{eq:des-split-parent-trace}
\end{equation}

Let $R_{\mathrm{con}}$ and $h_{\mathrm{con}}$ denote the projection and
homotopy in the tensor-product deformation retract of
Proposition~\ref{prop:des-contact-stokes}.  The same calculation takes place
before applying the trace: both monomials reduce to the same basis vector in
\eqref{eq:des-relative-normal-form}, with coefficients in the parenthesis in
\eqref{eq:des-split-parent-trace}.  Consequently
\begin{equation}
 R_{\mathrm{con}}
 (\beta^{\mathrm{sp}}_{J,\pi}+\beta^{\mathrm{par}}_{J,\pi},0,0,0)=0.
 \label{eq:des-split-parent-projection}
\end{equation}
This includes residue-three sectors because the relative normal form retains
all sixteen residues.

Let
\begin{equation}
 G^{\mathrm{con}}_{J,\pi}
 =h_{\mathrm{con}}(\beta^{\mathrm{sp}}_{J,\pi}
                  +\beta^{\mathrm{par}}_{J,\pi}),\qquad
 \delta G^{\mathrm{con}}_{J,\pi}
 =\beta^{\mathrm{sp}}_{J,\pi}
  +\beta^{\mathrm{par}}_{J,\pi}.
 \label{eq:des-corolla-filler}
\end{equation}
For the oriented interval $e=[v_0,v_1]$, define
\begin{equation}
 \kappa_{J,\pi}
 =e\otimes\beta^{\mathrm{sp}}_{J,\pi}
  +v_0\otimes G^{\mathrm{con}}_{J,\pi}.
 \label{eq:des-corolla-chain}
\end{equation}
For the tensor differential
$d(c\otimes\alpha)=\partial c\otimes\alpha
+(-1)^{\dim c}c\otimes\delta\alpha$, one has
\begin{equation}
 d\kappa_{J,\pi}
 =v_1\otimes\beta^{\mathrm{sp}}_{J,\pi}
  +v_0\otimes\beta^{\mathrm{par}}_{J,\pi}.
 \label{eq:des-corolla-boundary}
\end{equation}

Let $Q_F=[0,1]^{E_{\mathrm{int}}(F)}$.  The standard cellular contraction of
$Q_F$ to its all-zero vertex, tensored with the two-variable contraction
$(R_{\mathrm{con}},\iota,h_{\mathrm{con}})$ of
Proposition~\ref{prop:des-contact-stokes}, gives an operator $H_F$ with
\begin{equation}
 dH_F+H_Fd=1-\iota_F\Pi_F.
 \label{eq:des-forest-contraction}
\end{equation}
We now give the coherent extension, rather than choosing a primitive on each
square independently.  Order the nonempty faces of $Q_F$ first by dimension
and then lexicographically by their sets of free coordinates.  On vertices
and edges use the coefficients \eqref{eq:des-split-parent-seeds} and the corolla chain
\eqref{eq:des-corolla-chain}.  Suppose chains have been assigned to all
proper faces of a $k$-face $G$ and agree on their intersections.  Their
oriented sum is a cycle $Z_G$.  At the all-zero vertex, every refinement
order has the same block-sum exponent and product of vertex scalars.  On all
other terms \eqref{eq:des-split-parent-projection} gives zero normal-form
projection.  Thus
\begin{equation}
 dZ_G=0,\qquad \Pi_FZ_G=0,\qquad d(H_FZ_G)=Z_G.
 \label{eq:des-coherent-face-induction}
\end{equation}
Assign $H_FZ_G$ to $G$.  This recursion extends the chain across every face
and automatically uses the already fixed primitive on every proper face.  Because the order,
$H_F$, and the block-sum maps are fixed and equivariant, the construction
commutes with relabelling, coordinate transposition, restriction to a
subforest, and transport of the coefficient complex.  In particular, a
refinement cube receives one coherent extension, rather than unrelated
primitives on its squares; all higher compatibility follows from the induction in
\eqref{eq:des-coherent-face-induction}.

Define $C_*^{\mathrm{fac}}(F)$ to be the smallest differential submodule of
$C_*^{\mathrm{cell}}(Q_F)\otimes
\mathcal B_{\mathrm{con}}\otimes\mathfrak o_F$ containing every vertex coefficient,
every corolla chain \eqref{eq:des-corolla-chain}, and every recursively
assigned chain $H_FZ_G$, and closed under restriction to faces.  Identify a
length-one face with its cut forest through
\eqref{eq:des-forest-orientation}.  Its differential is
\begin{equation}
 d_{\mathrm{fac}}(c\otimes\alpha)
 =\partial_{\mathrm{cell}}c\otimes\alpha
  +(-1)^{\dim c}c\otimes\delta\alpha.
 \label{eq:des-factorization-differential}
\end{equation}
For a target frame $\Gamma$, transport the assembled coefficient into
$\mathcal B_{\Gamma,\mathrm{con}}$.  Define $\mu_\Gamma$ to be zero off
total degree zero and, on a zero-dimensional boundary cell, by
\begin{equation}
 \mu_\Gamma(c\otimes\alpha)
 =\operatorname{Tr}_\Gamma(\alpha).
 \label{eq:des-factorization-moment}
\end{equation}
Equations \eqref{eq:des-split-parent-trace},
\eqref{eq:des-corolla-boundary}, and
\eqref{eq:des-forest-contraction} prove
\begin{equation}
 \mu_\Gamma d_{\mathrm{fac}}=0.
 \label{eq:des-factorization-cocycle}
\end{equation}

\subsection{The cotangent-line zero cobordism}

Fix an exact open datum $\delta\in\mathcal E$, a distinguished label
$j_1\in I(\delta)$, and write
\begin{equation}
 \mathcal M_\delta=\mathcal{PM}_{\Gamma_\delta},\qquad
 E_\delta^-=\mathbb E_{\Gamma_\delta}(\mathbf d),\qquad
 E_\delta^+=\mathbb E_{\Gamma_\delta}(\mathbf d+\mathbf e_{j_1})
 \simeq E_\delta^-\oplus\mathcal L_{j_1}.
 \label{eq:des-open-bundles}
\end{equation}
The last isomorphism is the Witten-bundle splitting
\eqref{eq:intro-cotangent-splitting}.  Choose, simultaneously on the finite
boundary diagram of $\delta$, compatible equivariant rational PL
multisections $\sigma_\delta^-$ of $E_\delta^-$ and $s_\delta^+$ of
$\mathcal L_{j_1}$.  They are chosen by induction from the deepest boundary
strata, using relative PL extension on each finite proper-etale atlas.  Thus
their restrictions are the exterior sums prescribed by the bundle maps
$\phi_{\Lambda\Gamma}$, and no numerical Euler count is used in their
choice.

For every retained marking or boundary flag $w$, let
$t_{j_1,w}$ be the standard section of $\mathcal L_{j_1}$ obtained by
evaluating at $z_{j_1}$ the normalized meromorphic differential with residues
$+1$ at $w$ and $-1$ at the conjugate of $z_{j_1}$
\cite[Notation~6.1]{grossetal2022open}.  Relative PL extension gives a
transverse homotopy
\begin{equation}
 H_{\delta,w}:[0,1]\times\mathcal M_\delta
 \longrightarrow E_\delta^-\oplus\mathcal L_{j_1},
 \quad
 H_{\delta,w}(0)=\sigma_\delta^-\oplus t_{j_1,w},
 \quad
 H_{\delta,w}(1)=\sigma_\delta^-\oplus s_\delta^+.
 \label{eq:des-cotangent-homotopy}
\end{equation}
The construction is relative to every face where the sections are already
fixed.  On a boundary map that contracts the component carrying $j_1$, the
tautological morphism need not be invertible; that collar is retained as a
relative stop and is disjoint from the zero set by the rank and positivity
condition included in the exact datum.  On every other face,
\eqref{eq:des-cotangent-homotopy} pulls back under
$F_{\Lambda\Gamma}$ and $\phi_{\Lambda\Gamma}$.

Let $\mathcal Z_{\delta,w}$ be the determinant-oriented rational PL zero
chain of \eqref{eq:des-cotangent-homotopy}.  Its ordinary parameterized Euler
boundary is
\begin{equation}
 \partial\mathcal Z_{\delta,w}
 =[Z(\sigma_\delta^-\oplus s_\delta^+)]
 -[Z(\sigma_\delta^-\oplus t_{j_1,w})]
 +[Z(H_{\delta,w}|_{[0,1]\times\partial\mathcal M_\delta})].
 \label{eq:des-cotangent-zero-boundary}
\end{equation}
The first term is the Euler chain for descendent degree
$\mathbf d+\mathbf e_{j_1}$, by \eqref{eq:des-open-bundles}.  The divisor of
$t_{j_1,w}$ is the sum of open--closed boundary divisors on which $j_1$ lies
on the closed component and, when $w$ is an interior marking, $w$ remains on
the open component \cite[Lem.~6.4]{grossetal2022open}.  The local equation is
$u z^4$ in the quartic root coordinate, with $u$ a unit.  Intersecting with
$\sigma_\delta^-$ leaves exactly the rank-matched splittings $N_A=-1$; the
others remain as residual-degree faces with zero degree-zero pushforward.
Thus the degree shift and the open--closed term in the recursion are both
geometric terms of \eqref{eq:des-cotangent-zero-boundary}.

The signed coefficients introduced below do not weight branches of a
multisection.  They weight separate chains $\mathcal Z_{\delta,w}$.  Because
their sum is one, all copies of the high endpoint in
\eqref{eq:des-cotangent-zero-boundary} combine to the single canonical Euler
chain $Z(\sigma_\delta^-\oplus s_\delta^+)$.  Negative coefficients are
therefore legitimate rational-chain coefficients and never PL-multisection
weights.

\subsection{Compactified annular incidence spaces}

\subsubsection{Compactified incidence types}

The geometric source of the correspondence is a compactification of a
fiber product.  Its open stratum parametrizes a stable labelled tropical
forest in the annular corridor, together with an open quartic $W$-spin curve
at each root and a closed quartic $W$-spin curve at each detached component.
The evaluation positions satisfy the affine edge and wall-incidence
equations, while the Witten sections vanish on the corresponding moduli
factors.  The coefficient of a forest is the relative twisted de~Rham class
constructed above.

The compactification records each geometric degeneration rather than
identifying it numerically.  Zero and infinite forest lengths give
contractions and cuts; wall crossings acquire their Cartan homotopies;
colliding directions are resolved by Fulton--MacPherson spaces; the universal
boundary circle records the pointing direction; and fourth-root collars
record the normal direction to an open--closed splitting.  The usual compact
open and closed $W$-spin spaces supply their own boundary strata.  Auxiliary
length caps are placed beyond the incidence locus.  These constructions give
the following local model.

\begin{definition}[Stable resolved type and incidence section]
\label{def:des-resolved-type}
A \emph{stable resolved labelled type} $\tau$ consists of the following
finite data.
\begin{enumerate}
\item It has a rooted forest with leaf set contained in $I$.  Every internal
vertex is either at least trivalent or carries one of the following specified
geometric events: crossing a retained wall; identifying two affine-chart
copies on an overlap; contracting or cutting an internal forest edge;
colliding two tangent directions; moving the pointing flag on the compactified
boundary circle; or splitting off a closed $W$-spin component.  In
particular, unmarked bivalent wall-free vertices are suppressed.  Each open
root component names an exact datum $\delta$ of
\eqref{eq:intro-exact-open-datum}, including its boundary maps and Witten
bundles.
\item Every oriented segment $e$ has a fixed integral slope $v_e$, a
nonnegative length $\lambda_e$, and an ordered word in the finitely many
walls retained modulo $\mathfrak m_S^N$.  Each nonidentity bend has positive
deformation order, and the sum of the orders on a path is less than $N$.
\item Every detached component names its compact quartic $W$-spin base
$\mathcal{PM}_{\delta}$ and Euler bundle $\mathbb E_{\delta}$ from
$\mathcal E$.  Every coincident set of directions names the corresponding
Fulton--MacPherson factor.  At a vertex carrying $j_1$, the oriented real
blowup of its tangent-direction configuration is the compactified universal
boundary circle.  Its named vertices are the root flag, retained markings,
the two Fermat phase flags, wall directions, collision directions, and
closed-component directions.  A type chooses one closed cell $M_\tau$ of
this circle and one cotangent-homotopy interval
$[0,1]_{\psi,\delta,w}$ from \eqref{eq:des-cotangent-homotopy}; without a
distinguished label these factors are points.  Adjacent cells are identified
at their common endpoints in the face category below.
\item The type contains a nondegenerate composable chain of $q$ boundary or
overlap maps from \eqref{eq:intro-exact-open-datum}, and an ordered word of
$b$ retained wall maps.  Boundary arrows strictly lower the graph-complexity
order of Appendix~\ref{app:gkt-completion}.  For overlaps, fix a finite
ordered corridor cover and retain only the strictly increasing \v{C}ech chains of
distinct chart indices.  The inverse transition is the coefficient map on
the same overlap, not a second nerve arrow; identities and repeated indices
are omitted.  Mixed chains use the product normalization: a boundary arrow
retains the ordered overlap multi-index, and an overlap arrow retains the
graph type.  The corresponding simplices
$\Delta^q_{\mathrm{nerve}}$ and $\Delta^b_{\mathrm{bar}}$ record composition
of the boundary maps and the ordered Cartan homotopy, respectively.  Every
affine overlap constraint records as part of $\tau$ the restriction
$(A_j,c_j)$ of the fixed integral-affine transition
$x^+=A_jx^-+c_j$ between its two corridor cells.
\end{enumerate}

Let $\mathcal V_\tau$ be the position vertices, $\mathcal S_\tau$ the
segments, $\mathcal H_\tau$ the affine-overlap parameters,
$\mathcal E_\tau$ the internal forest edges, $\mathcal W_\tau$ the prescribed
wall incidences, $\mathcal M_\tau$ the overlap and common-vertex matching
constraints, and $\mathcal D_\tau$ the closed-component smoothing roots.  Let
$\mathcal E_\tau^{\mathrm{op}}$ and
$\mathcal E_\tau^{\mathrm{cl}}$ be the named open and closed moduli data.
$\mathcal Q_\tau$ is the set of chosen pairs $(\delta,w)$ carrying a
cotangent homotopy.
A constraint $j\in\mathcal M_\tau$ has
dimension $d_j\in\{1,2\}$.  In a chamber $\mathfrak c$, put
\begin{equation}
 A_{\tau,\mathfrak c}=
 \left\{(x_v,\lambda_e):
 \begin{array}{ll}
 x_v\in\RR^2 &(v\in\mathcal V_\tau),\\
 0\leq\lambda_e\leq L_e &(e\in\mathcal S_\tau),\\
 \epsilon_{\mathfrak c h}\ell_h(x_{v(h)})\geq0
   &(h\text{ a corridor side})
 \end{array}\right\}.
 \label{eq:des-position-polytope}
\end{equation}
The cap $L_e$ lies beyond the incidence zero set.  For the specified chains
of lengths $q$ and $b$, define
\begin{equation}
\begin{aligned}
 P_\tau={}&A_{\tau,\mathfrak c}
 \times[0,1]^{\mathcal H_\tau}
 \times[0,1]^{\mathcal E_\tau}\\
 &\times\Delta^q_{\mathrm{nerve}}
 \times\Delta^b_{\mathrm{bar}}
 \times\prod_C\operatorname{FM}(C)\\
 &\times M_\tau
 \times\prod_{(\delta,w)\in\mathcal Q_\tau}[0,1]_{\psi,\delta,w}
 \times\prod_{D\in\mathcal D_\tau}
   \bigl([\mathbb D_R/\mu_4]\times[0,1]_D\bigr)
 \times\prod_{\delta\in\mathcal E_\tau^{\mathrm{op}}}
   \mathcal M_\delta
 \times\prod_{\delta\in\mathcal E_\tau^{\mathrm{cl}}}
   \mathcal{PM}_{\delta}.
\end{aligned}
 \label{eq:des-ambient-product}
\end{equation}
Each factor in
\eqref{eq:des-ambient-product} occurs once.

The oriented equation orbibundle $E_\tau\to P_\tau$ and its section
$s_\tau$ are the following ordered direct sums:
\begin{equation}
\begin{array}{c|c|c}
\text{index}&\text{bundle summand}&\text{component of }s_\tau\\ \hline
e\in\mathcal S_\tau&\RR^2&x_{h(e)}-x_{t(e)}-\lambda_ev_e\\
h\in\mathcal W_\tau&\RR&\ell_h(x_{v(h)})\\
j\in\mathcal M_\tau&\RR^{d_j}&x_j^+-A_jx_j^--c_j\\
D\in\mathcal D_\tau&\CC&z_D^4-t_D\\
\delta\in\mathcal E_\tau^{\mathrm{op}}&
 E_\delta^-\oplus\mathcal L_{j_1}&H_{\delta,w}(u_\psi)\\
\delta\in\mathcal E_\tau^{\mathrm{cl}}&
 \mathbb E_\delta&s_\delta^{\mathrm{Euler}}.
\end{array}
 \label{eq:des-incidence-section}
\end{equation}
For a common-vertex constraint $(A_j,c_j)=(1,0)$; for an overlap constraint
it is the integral-affine transition recorded in Item~(4).  The open-bundle
component is
\eqref{eq:des-cotangent-homotopy}, so its $u_\psi=1$ restriction is the
degree-$\mathbf d+\mathbf e_{j_1}$ Euler section and its $u_\psi=0$
restriction contains the pointing divisor.  Thus every component map and
its real rank are specified.  Put
\begin{equation}
\begin{aligned}
 d_\tau^{\mathrm{op}}
 &=\sum_{\delta\in\mathcal E_\tau^{\mathrm{op}}}
   \dim_{\RR}\mathcal M_\delta,\\
 d_\tau^{\mathrm{cl}}
 &=\sum_{\delta\in\mathcal E_\tau^{\mathrm{cl}}}
   \dim_{\RR}\mathcal{PM}_\delta,\\
 r_\tau^{\mathrm{mod}}
 &=\sum_{\delta\in\mathcal E_\tau^{\mathrm{op}}}
   \operatorname{rank}_{\RR}(E_\delta^-\oplus\mathcal L_{j_1})
   +\sum_{\delta\in\mathcal E_\tau^{\mathrm{cl}}}
   \operatorname{rank}_{\RR}\mathbb E_\delta.
\end{aligned}
 \label{eq:des-moduli-dimensions}
\end{equation}
In particular,
\begin{equation}
\begin{aligned}
 \operatorname{vdim}(\tau)={}&
 2|\mathcal V_\tau|+|\mathcal S_\tau|+|\mathcal H_\tau|
 +|\mathcal E_\tau|+q+b+\sum_C\dim\operatorname{FM}(C)
 +\dim M_\tau+|\mathcal Q_\tau|\\
 &+3|\mathcal D_\tau|
 +d_\tau^{\mathrm{op}}+d_\tau^{\mathrm{cl}}
 -2|\mathcal S_\tau|-|\mathcal W_\tau|-\sum_jd_j
 -2|\mathcal D_\tau|
 -r_\tau^{\mathrm{mod}},
 \label{eq:des-virtual-dimension}
\end{aligned}
\end{equation}
and its orientation line is
\begin{equation}
 \mathfrak o_\tau=\det TP_\tau\otimes(\det E_\tau)^\vee.
 \label{eq:des-incidence-orientation}
\end{equation}
All subsequent chain identities are degreewise in this virtual grading.
\end{definition}

There are only finitely many stable types.  Indeed, a stable forest on a
fixed finite leaf set has finitely many unmarked vertices; every remaining
vertex carries one of the finite named events, and a path has at most $N-1$
positive-order bends.  The finite wall list, finite dependency closure,
finite profile sequence \eqref{eq:intro-balanced-profiles}, and finite
subdivisions then leave a finite set of choices.  Moreover, the length $q$
is bounded by the sum of the number of graph-complexity levels and the number
of charts in the ordered corridor cover, because every nonidentity step
strictly advances one of these two finite ranks.  Hence the normalized
\v{C}ech--boundary chains also leave only finitely many choices.

\begin{proposition}[Boundary stratification]
\label{prop:des-boundary-stratification}
Every real codimension-one stratum of the compactified incidence space is one
of the following: an annular endpoint; a contraction or cut of a tropical
edge; a wall, joint, or affine-overlap stratum; a collision of tangent
directions; an endpoint of the cotangent homotopy or the fourth-root normal
collar; or a boundary stratum of an open or closed $W$-spin moduli factor.
The auxiliary caps and disk boundaries are disjoint from the perturbed zero
locus.  After the canonical identifications on paired faces, the exterior
boundary consists precisely of the canonical cotangent face, open--closed
splittings, wall crossings, cotangent contractions, pairwise exchanges, and
closed factors of nonzero Euler degree, together with the two annular
restrictions.
\end{proposition}

\begin{proof}
Every facet $D\subset P_\tau$ is obtained by placing one factor in
\eqref{eq:des-ambient-product} on a codimension-one face.  It has a resolved
type $\tau_D$ and a map $f_D:P_{\tau_D}\to D$.  The following table records
the local chart map and its normal-bundle comparison.  A ``trivial summand''
is a coordinate together with a constraint on which the displayed
elimination is an oriented isomorphism.
\begin{center}
\begin{tabularx}{\textwidth}{@{}
 >{\raggedright\arraybackslash}p{.21\textwidth}
 >{\raggedright\arraybackslash}p{.25\textwidth}
 >{\raggedright\arraybackslash}X@{}}
\toprule
Facet of $P_\tau$ & Target type $\tau_D$ & Map and normal-bundle comparison\\
\midrule
$\lambda_e=0$ or an overlap-cylinder end & contracted or stabilized tree & substitute $x_{h(e)}=x_{t(e)}$ or $x^+=A_jx^-+c_j$ and delete the corresponding trivial summand\\
$\lambda_e=L_e$ or a corridor end & empty, respectively $\partial_{\mathrm{ann}}^\pm$ & the auxiliary cap is disjoint from the zero set; corridor ends are retained for the restriction maps $e^\pm$\\
forest metric $t_e=0,1$ & corolla, respectively cut forest & at $0$ use \eqref{eq:des-corolla-chain}; at $1$ use $\Phi_F$ and delete the matching trivial summand\\
wall-side equality & wall, joint, tangency, or bend-collision type & conjugate the affine constraints by the labelled transition; use the bracket square at a joint\\
boundary-map simplex face & first/last map or composite of adjacent maps & restrict, delete, or compose the corresponding geometric maps\\
ordered wall-homotopy face & shorter wall word or bracket joint & apply the chain map $\mathsf T_H$ or the homotopy $\mathsf P_H$ defined below\\
collision scale $0$ & exchanged pair or collision block & restrict to the real blowup face and split off its oriented normal\\
moving-circle vertex & root, pointing, contraction, wall, or closed-splitting direction & restrict the same open bundle and pointing section to the named vertex of $M_\tau$\\
cotangent parameter $u_\psi=0,1$ & pointing divisor, respectively degree-$\mathbf d+\mathbf e_{j_1}$ Euler chain & use the two endpoint sections in \eqref{eq:des-cotangent-homotopy}\\
root parameter $t_D=0,1$ & closed splitting, respectively pointing face & restrict $z_D^4-t_D$ to the indicated endpoint\\
open-moduli boundary & the named boundary graph $\Lambda$ & apply $(F_{\Lambda\Gamma},\phi_{\Lambda\Gamma})$ from \eqref{eq:intro-exact-open-datum}\\
closed degeneration & the named dependent closed datum in $\mathcal E$ & apply its stated boundary map and restrict the single Euler bundle already present\\
\bottomrule
\end{tabularx}
\end{center}
Auxiliary caps and disk boundaries are disjoint from the zero set.  The
closed orbifold bases and moving circles have no other real
codimension-one faces.

The stable cancellations used in these local comparisons are
\begin{equation}
\begin{array}{c|c|c}
\text{facet}&\text{substitution}&\text{eliminated coordinate and constraint}\\ \hline
\lambda_e=0&x_{h(e)}=x_{t(e)}&(x_{h(e)},\ x_{h(e)}-x_{t(e)})\\
\text{overlap end}&x^+=A_jx^-+c_j&(x^+,\ x^+-A_jx^--c_j)\\
t_e=1&\text{cut the forest at }e&(t_e,\text{ matching constraint})\\
z_D^4=t_D& t_D=z_D^4&(t_D,\ z_D^4-t_D).
\end{array}
\label{eq:des-row-eliminations}
\end{equation}
The determinant of each displayed coordinate-to-constraint derivative is
positive in the listed order.  At $t_e=0$, the coefficient constraint is not removed: it is the
boundary of \eqref{eq:des-corolla-chain}.  At an open or closed moduli face,
the oriented bundle comparison is the explicitly included map
$\phi_{\Lambda\Gamma}$.

Restricting the ordered components of \eqref{eq:des-incidence-section} along a
facet gives a stable oriented identification
\begin{equation}
 f_D^*(E_\tau|_D)\oplus\underline{\RR}^{\,a_D}
 \simeq E_{\tau_D}\oplus\underline{\RR}^{\,b_D},
 \label{eq:des-stable-face-bundle}
\end{equation}
where the trivial summands are precisely the stable cancellations in
\eqref{eq:des-row-eliminations}.  The maps around a codimension-two face are
equal because affine substitution is associative, the exact-data maps obey
their actual composition identities, and the coherent forest primitives are
defined on the whole refinement cube by
\eqref{eq:des-coherent-face-induction}.  Taking
determinants removes one outward normal.  Interchanging two facet normals
reverses their exterior product and proves the two-face boundary sign even
when the rank of $ds_\tau$ changes.
At a partition face, permutation of incoming affine blocks acts by the sign
character on \eqref{eq:des-incidence-orientation}; the same character acts
on $\mathfrak o_F$ in \eqref{eq:des-forest-orientation}.  Their tensor is
therefore relabelling invariant.
The displayed product has no further kinds of codimension-one face, which
proves the asserted exhaustion.
\end{proof}

\subsubsection{Virtual PL zero loci}

Choose simultaneous equivariant triangulations of the finite proper
etale atlases, the bundles, the sections, and all face maps.  After relative
simplicial approximation, $s_\tau$ is PL in each stable bundle chart.  A
rational PL multisection on a chart is a finite set of PL branches with
positive rational weights summing to one.  On overlaps, replace two branch
sets by their finite weighted common refinement; at isotropy include the
complete finite orbit with equal subdivided weights, and at product corners
take Cartesian products of branch sets.

These positive numbers are used only to perturb the incidence section.
They are distinct from the signed rational coefficients of the moving-circle
chains used in the cotangent recursions below.

Induct on face depth and simplex dimension.  Prescribed branches on a
downward-closed family of faces extend across each new simplex, and their
new vertex values can be chosen outside the finite union of affine
subspaces for which a branch fails PL general position.  The stable
identifications \eqref{eq:des-stable-face-bundle} make the choices agree
across every face and orbifold chart.  The same induction on
$[0,1]\times P_\tau$ joins any two systems relative to the prescribed faces.
Thus compatible arbitrarily small multisections $\nu_\tau$ exist and
\begin{equation}
 [Z_\tau]^{\mathrm{vir}}=[(s_\tau+\nu_\tau)^{-1}(0)]
 \label{eq:des-virtual-zero}
\end{equation}
is a compact determinant-oriented rational branched PL chain of dimension
$\operatorname{vdim}(\tau)$ whose boundary is its intersection with the
ambient facets.

\subsubsection{Wall, joint, and gluing homotopies}

For completeness, the coefficient homotopies used on wall facets are fixed
as follows.  If $H$ is a retained wall Hamiltonian, contraction with its
Hamiltonian vector field gives a degree-one operator $\mathsf K_H$ on the
transported contact complex.  Put
\begin{equation}
 \boldsymbol\Delta_H=[\delta,\mathsf K_H],\qquad
 \mathsf T_H=\exp(\boldsymbol\Delta_H),\qquad
 \mathsf P_H=\sum_{n=0}^{N-1}
   \frac{\mathsf K_H\boldsymbol\Delta_H^n}{(n+1)!}.
 \label{eq:des-explicit-wall-homotopy}
\end{equation}
All sums are finite modulo $\mathfrak m_S^N$, and the Cartan commutator
identity gives
\begin{equation}
 \delta\mathsf P_H+\mathsf P_H\delta
 =\mathsf T_H-\operatorname{id}.
 \label{eq:des-wall-homotopy}
\end{equation}
For an ordered word $H_1\cdots H_b$, use the telescoping homotopy
\begin{equation}
 \mathsf P_{H_b\cdots H_1}
 =\sum_{i=1}^b
  \mathsf T_{H_b}\cdots\mathsf T_{H_{i+1}}
  \mathsf P_{H_i}.
 \label{eq:des-word-homotopy}
\end{equation}
At a two-wall joint the two boundary words have equal transport by scattering
consistency.  Let $Z$ be the difference of their two instances of
\eqref{eq:des-word-homotopy}.  It is a degree $-1$ cycle in the endomorphism
complex of the contracted contact complex.  The retract
\eqref{eq:des-relative-normal-form} is concentrated in cohomological degree
$2$, while the contact complex has no terms above degree $2$; hence
$R_{\mathrm{con}}Z=0$.  The contraction identity now gives
$[\delta,h_{\mathrm{con}}Z]=Z$, so $h_{\mathrm{con}}Z$ is the specified
square primitive.  At a joint of codimension $r\geq3$, the oriented sum of
the already assigned proper faces is a cycle of degree $1-r$.  Its projection
to the normal form again vanishes by degree, and the same contraction fills
it.  Induction over the ordered faces, exactly as in
\eqref{eq:des-coherent-face-induction}, therefore supplies a compatible
cubical primitive.  This defines the source, target, and signs of every wall
and joint map used below.

\subsubsection{The cellular coend}

Choose the simultaneous subdivisions so that every map in the
boundary-stratification table is
cellular.  Let $\mathcal K_{I,\mathbf d,N}$ be their finite face category.
Its objects are pairs $(\tau,\sigma)$, and a generating arrow
\begin{equation}
 f:(\tau_D,\sigma_D)\longrightarrow(\tau,\sigma)
 \label{eq:des-face-arrow}
\end{equation}
is either the inclusion of an ordinary codimension-one subdivision face
$\sigma_D\subset\sigma$ with $\tau_D=\tau$, or the inclusion of a cell lying
in one of the ambient facets $D$ in the table followed by its stated map.
The coefficient map for the first kind is the identity.  Composites are the
actual composites of restrictions, exact-data maps, and the fixed primitives
above.  Hence the two arrows obtained from a codimension-two corner are equal,
rather than imposed by an unnamed relation.  Dimension strictly increases
along a nonidentity arrow, so the category is direct; its finiteness follows
from the finiteness proved after Definition~\ref{def:des-resolved-type}.

Let $\mathcal G_{\tau,\sigma}$ be the rational PL branched-chain group on
$Z_\tau^{\mathrm{vir}}\cap\sigma$, relative to its intersections with the two
corridor-end facets $\partial_{\mathrm{ann}}^-\cup
\partial_{\mathrm{ann}}^+$.  Collar restriction to either relative facet is
retained as a separate map $e^\pm$ below.  In the transported contact complex,
let
$\mathcal B^{\mathrm{sel}}_{\tau,\sigma}$ be the smallest differential
submodule containing the selected coefficient $\alpha_\tau$, all restrictions of
the coherent forest primitives, all operators
\eqref{eq:des-explicit-wall-homotopy}--\eqref{eq:des-word-homotopy} applied to
those elements, and all their root, open-boundary, and closed-boundary images.
It is finite at the chosen deformation order.  If
$\Gamma(\tau,\sigma)$ is the chamber frame, set
\begin{equation}
 \mathcal C_{\tau,\sigma}
 =\mathcal B^{\mathrm{sel}}_{\tau,\sigma}
  \otimes\mathfrak o_{F_\tau}.
 \label{eq:des-cell-coefficient-complex}
\end{equation}
This is the differential subcomplex generated by every coefficient and
primitive used in the construction, not the line spanned by one selected
coefficient.
The arrow \eqref{eq:des-face-arrow} induces
$f_*:\mathcal G_{\tau_D,\sigma_D}\to\mathcal G_{\tau,\sigma}$ and a chain
map
\begin{equation}
 f^\#:\mathcal C_{\tau,\sigma}\longrightarrow
       \mathcal C_{\tau_D,\sigma_D}.
 \label{eq:des-coefficient-pullback}
\end{equation}
For an ordinary subdivision face it is the identity.  For an ambient facet it
is restriction followed, according to the table, by chamber transport, the
determinant map $\Phi_F$, the exact-data bundle map
$\phi_{\Lambda\Gamma}$, or the stated root or closed-boundary map.  On a
metric-zero generator it uses $G^{\mathrm{con}}$ from
\eqref{eq:des-corolla-filler}; on a metric-one generator it uses the cut-forest
element.  The source module was closed under exactly these images, so the map
is defined on every generator and extends linearly.  Equations
\eqref{eq:des-trace-transport}, \eqref{eq:des-corolla-boundary},
\eqref{eq:des-wall-homotopy}, and the exact-data composition identities prove
$f^\#\delta=\delta f^\#$ and equality on every composite.

Define
\begin{equation}
 \mathsf C^{\mathrm{tot}}_{I,\mathbf d,N}
 =\left(\bigoplus_{(\tau,\sigma)\in\mathcal K_{I,\mathbf d,N}}
 \mathcal G_{\tau,\sigma}\otimes\mathcal C_{\tau,\sigma}\right)
 \Big/\langle[f_*z,\alpha]-[z,f^\#\alpha]\rangle_f.
 \label{eq:des-coend}
\end{equation}
The closed factor, metric cube, boundary-map simplex, and ordered
wall-homotopy simplex occur in the branch polyhedron once; they are not
repeated in the coefficient.  The differential
is
\begin{equation}
 d[z,\alpha]=[\partial_{\mathcal K}z,\alpha]
 +(-1)^{\dim z}[z,\delta\alpha],
 \label{eq:des-total-differential}
\end{equation}
where $\partial_{\mathcal K}$ is the ordinary outward-normal-first cellular
boundary followed by the corresponding subdivision or ambient-facet arrow
\eqref{eq:des-face-arrow}.  This
is an endomorphism because \eqref{eq:des-cell-coefficient-complex} is closed
under $\delta$ and \eqref{eq:des-coefficient-pullback} is a chain map.  The
ordinary two-face cancellation gives $\partial_{\mathcal K}^2=0$; the Koszul
sign cancels the mixed terms.  Hence $d^2=0$ on the displayed quotient.

\subsubsection{Root thimbles and exterior faces}

The remaining root-normal collar is explicit.  The four radial branches of
$z^4=t$, $0\leq t\leq1$, descend to a chain $\mathfrak T_4$ with
\begin{equation}
 \partial[\mathfrak T_4]=[\mathrm{pt}]-4[B\mu _4].
 \label{eq:des-thimble}
\end{equation}
Since $[B\mu _4]$ pushes forward to $\frac14[\mathrm{pt}]$, its normalized
low endpoint has coefficient one.  On the already present closed and
continuing annular factors, this gives a thimble with boundary
$\partial_{\mathrm{pt}}-\partial_{\mathrm{oc}}$.

The wall and joint cylinders are the explicit Cartan homotopies
\eqref{eq:des-explicit-wall-homotopy}--\eqref{eq:des-word-homotopy} and their
coherent square primitives.  A multiple-collision cone is the radial cone over
the corresponding Fulton--MacPherson exceptional face; its bottom has
smaller dimension and hence zero cellular pushforward.

The named boundary chains are the following oriented sums:
\begin{center}
\begin{tabularx}{\textwidth}{@{}>{\raggedright\arraybackslash}p{.20\textwidth}X@{}}
\toprule
Chain & Defining codimension-one faces \\
\midrule
$\partial_{\mathrm{can}}$ & the $u_\psi=1$ zero chain of $\sigma_\delta^-\oplus s_\delta^+$, representing descendent degree $\mathbf d+\mathbf e_{j_1}$ \\
$\partial_{\mathrm{pt}}$ & the $u_\psi=0$ pointing-section divisor before the radial root collar is added \\
$\partial_{\mathrm{oc}}$ & the low root-thimble endpoint, multiplied by the compact closed factor and the continuing open factor \\
$\partial_{\mathrm{wall}}$ & critical open-moduli boundary graphs carrying the Cartan wall homotopy \\
$\partial_{\psi}$ & contraction boundary graphs of the cotangent homotopy, weighted by the one- or two-marking pointing chain \\
$\partial_{\mathrm{exch}}$ & the pair of cyclic orders on a two-direction collision face \\
$\partial_{\mathrm{deg}}$ & closed-factor faces whose Euler degree is incompatible with the required zero-dimensional pushforward \\
\bottomrule
\end{tabularx}
\end{center}
Metric-zero faces are filled by
\eqref{eq:des-corolla-chain}, metric-one faces are glued to cut forests,
structural-nerve faces cancel through the coend, wall and joint faces are
filled by \eqref{eq:des-word-homotopy} and its coherent square version, and
multiple collisions are filled by their radial cones.  Auxiliary caps and
root-chart boundaries are disjoint from the zero locus for all sufficiently
small perturbations: compactness gives a positive lower bound for
$\lVert s_\tau\rVert$ there.  These are all facets because
\eqref{eq:des-ambient-product} is a finite product of the half-space
polytope \eqref{eq:des-position-polytope}, cubes, simplices,
Fulton--MacPherson spaces, cotangent and root intervals, and the compact open
and closed orbifold factors in \eqref{eq:des-ambient-product}.

\subsubsection{Assembly and the boundary theorem}

For each resolved type, let $\alpha_\tau$ be the transported coefficient
assembled from its decorated tropical forest by
\eqref{eq:des-single-seed}.  A
\emph{compatible moving-circle chain} $\mathfrak p$ assigns a signed rational
coefficient $w_\tau(\mathfrak p)$ to every oriented cell $M_\tau$, with equal
coefficients on identified subdivisions.  With product orientation ordered
as in \eqref{eq:des-ambient-product}, define
\begin{equation}
 \Psi^{\mathrm{raw}}_{I,\mathbf d,N}[\mathfrak p]
 =\sum_\tau w_\tau(\mathfrak p)
   [Z_\tau^{\mathrm{vir}},\alpha_\tau].
 \label{eq:des-raw-chain}
\end{equation}
Here the open-bundle component of $Z_\tau^{\mathrm{vir}}$ is the
parameterized zero chain $\mathcal Z_{\delta,w}$ of
\eqref{eq:des-cotangent-zero-boundary}.  Hence its named outward boundary is
$\partial_{\mathrm{can}}-\partial_{\mathrm{pt}}
+\partial_{\mathrm{wall}}+\partial_{\psi}
+\partial_{\mathrm{exch}}+\partial_{\mathrm{deg}}$; its remaining faces are precisely
the internal faces listed above.  Add to \eqref{eq:des-raw-chain} the wall
and joint prisms, corolla and refinement primitives, collision cones, and
radial thimbles, each with the sign opposite to its incident internal face.
Denote the resulting finite sum by $\Psi_{I,\mathbf d,N}[\mathfrak p]$.

\begin{theorem}[Cotangent-homotopy boundary]
\label{thm:des-absolute-boundary}
For every compatible moving-circle chain $\mathfrak p$, the chain
$\Psi_{I,\mathbf d,N}[\mathfrak p]$ is well defined in
\eqref{eq:des-coend} and satisfies
\begin{equation}
 \partial\Psi_{I,\mathbf d,N}[\mathfrak p]
 =\partial_{\mathrm{can}}-\partial_{\mathrm{oc}}
  +\partial_{\mathrm{wall}}+\partial_{\psi}
  +\partial_{\mathrm{exch}}+\partial_{\mathrm{deg}}.
 \label{eq:des-absolute-boundary}
\end{equation}
Its relative PL bordism class is independent of all triangulations,
multisections, auxiliary primitives, and chosen root arcs, and is natural under window
restriction, deformation-order reduction, finite-root refinement,
relabelling, target-frame re-expansion, coordinate transposition, and
common-enlargement transport by $K$.
\end{theorem}

\begin{proof}
Equation \eqref{eq:des-cotangent-zero-boundary} supplies
$\partial_{\mathrm{can}}-\partial_{\mathrm{pt}}$ and the open-moduli side
faces.  The ordinary cellular boundary of $\mathfrak p$ gives their signed
pointing and contraction coefficients.  The chain-homotopy identity
\eqref{eq:des-wall-homotopy} cancels each wall
face, and its two-parameter form cancels each joint face.  Equations
\eqref{eq:des-corolla-boundary} and
\eqref{eq:des-factorization-cocycle} fill zero-length forest faces and glue
length-one faces.  The two cyclic orders at a pair collision have opposite
outward-normal orientations and form the displayed exchange pair.  For a collision
of at least three directions, the bottom of the radial cone has smaller
dimension and zero cellular pushforward; nested cone sides cancel in pairs.
Simplicial boundary-map, affine-overlap, and stabilization faces cancel by the defining
coend relation.  Coefficient-order terms in $\mathfrak m_S^N$ vanish, while
closed degree mismatches are retained in $\partial_{\mathrm{deg}}$.  The
corridor-end faces are relative and are recorded by $e^\pm$.  Finally,
\eqref{eq:des-thimble} changes the raw boundary
$\partial_{\mathrm{can}}-\partial_{\mathrm{pt}}$ into
$\partial_{\mathrm{can}}-\partial_{\mathrm{oc}}$, with coefficient
$4\cdot\frac14=1$.  The product-facet classification preceding
\eqref{eq:des-raw-chain} proves that no other codimension-one face remains.

For window restriction, deformation reduction, or finite-root refinement,
the smaller face diagram is a downward-closed subdiagram of the larger one.
Choose the multisections by relative extension along a nested cofinal
exhaustion.  Relabelling signs cancel between
\eqref{eq:des-incidence-orientation} and
\eqref{eq:des-forest-orientation}; coordinate transposition is obtained by
taking the two-element branch orbit.  Transport to another target frame and
core transport after common enlargement carry the incidence equations, the
open Witten-bundle diagram, the relative contact complex, and every chosen primitive by
chain isomorphisms.  These maps therefore
commute with \eqref{eq:des-total-differential}.  The parameterized
multisection construction gives the claimed compatible cobordisms and
independence of the exhaustion.
\end{proof}

\subsection{Cotangent recursions}

Let $i^\pm:C^\pm\hookrightarrow C$ be the incoming and outgoing facets of
the compact corridor.  Restriction of branch polyhedra along $i^\pm$ and the
coefficient pullback \eqref{eq:des-coefficient-pullback} define chain maps
\begin{equation}
 e^\pm:\mathsf C^{\mathrm{tot}}_{I,\mathbf d,N}
 \longrightarrow\mathsf C^{\mathrm{end},\pm}_{I,\mathbf d,N}.
 \label{eq:des-endpoint-maps}
\end{equation}
Here $\mathsf C^{\mathrm{end},\pm}_{I,\mathbf d,N}$ is the same coend
\eqref{eq:des-coend} over the full subcategory of cells in
$\partial_{\mathrm{ann}}^\pm$; hence its objects, differential, and
coefficient maps are explicit restrictions of those already defined.  Both
targets are transported to the fixed frame $\Gamma$ using
\eqref{eq:des-transported-contact-complex}.

For any descendent vector $\mathbf d$, let
$\mathcal E^\pm_{I,\mathbf d,N}$ be the restriction to the corresponding
annular end of the canonical Euler zero chain
$Z(\mathbb E_\Gamma(\mathbf d),\sigma_\Gamma(\mathbf d))$.  Relative
extension makes these choices compatible in $\mathbf d$ and in the finite
window.  In particular, the canonical cotangent-homotopy boundary satisfies
\begin{equation}
 e^\pm(\partial_{\mathrm{can}}
       \Psi_{I,\mathbf d,N}[\mathfrak p])
 =\mathcal E^\pm_{I,\mathbf d+\mathbf e_{j_1},N}.
 \label{eq:des-endpoint-degree-shift}
\end{equation}
The moving-circle primitive has zero collar restriction, so the right side
does not depend on that primitive.  Define the annular moment by
\begin{equation}
 \mathcal A_{\mathrm{ann}}(I,\mathbf d)
 =\mu_\Gamma(\mathcal E^-_{I,\mathbf d,N}).
 \label{eq:des-annular-moment}
\end{equation}
Wall transport is a chain map, $\mu_\Gamma$ is a cocycle, and the intervening
relative chain has boundary \eqref{eq:des-absolute-boundary}; the wall,
exchange, and residual-degree terms have zero contact trace as proved below.
Thus the greatest-slope endpoint gives the same scalar after transport to the
same target frame.  The empty endpoint is the unit:
$\mathcal A_{\mathrm{ann}}(\varnothing,0)=1$.

For $A\subseteq I$, write
\begin{equation}
 \sum_{i\in A}a_i=r_A+4\ell_A,\qquad
 \sum_{i\in A}b_i=s_A+4m_A,
 \qquad
 N_A=\ell_A+m_A-|A|+1+\sum_{i\in A}d_i,
 \label{eq:des-threshold}
\end{equation}
with $0\leq r_A,s_A\leq3$.  When $N_A=-1$, let $z_A$ be the
degree-zero open half-node of twist $(r_A,s_A)$ and set
$Q_A=(I\setminus A)\sqcup\{z_A\}$.  Its compact closed partner has the
complementary twist $(2-r_A,2-s_A)$, interpreted modulo $4$.  Denote its
Frobenius-normalized Euler pushforward, including the degrees
$\mathbf d|_A$, by
\begin{equation}
 \mathcal C_A^{\mathrm{ext}}
 =16\int_{[\mathfrak C_A]}e(\mathbb E_A^{\mathrm{cl}}).
 \label{eq:des-closed-factor}
\end{equation}
The factor \eqref{eq:des-closed-factor} is an input from the compact closed
quartic $W$-spin moduli problem, not a value of an open invariant.

The wall-crossing trace vanishes by the primitive
\begin{equation}
 \mathcal K_{u,v}=(u+1)X^uY^{v+1}dX+(v+1)X^{u+1}Y^vdY,
 \label{eq:des-wc-primitive}
\end{equation}
whose differential is
\begin{equation}
 D_{W_0}\mathcal K_{u,v}
 =4(v+1)X^{u+4}Y^v\Omega-4(u+1)X^uY^{v+4}\Omega.
 \label{eq:des-wc-boundary}
\end{equation}
Indeed, the two terms in \eqref{eq:des-wc-boundary} have traces
$4(v+1)w(u+4,v)$ and $4(u+1)w(u,v+4)$, which agree by
\eqref{eq:des-gamma-weight}.  The complete relative differential also
contains its divisor restrictions, so Proposition~\ref{prop:des-contact-stokes}
annihilates every wall-crossing face.  The exchange trace vanishes because exchange of the
two directions reverses the real face orientation and preserves its
coefficient.

\subsubsection{Pointing chains}

We next define the two moving-circle chains and compute their contributions.
A rooted critical decomposition is
\[
 Q=((I_0,u,v);(I_1,p_1,q_1),\ldots,(I_h,p_h,q_h)),
\]
where $j_1\in I_0$, the root profile $(u,v)$ is critical, and the other
profiles are balanced.  Its two quartic phase faces and its $j$th
contraction face are
\begin{equation}
\begin{aligned}
 P_X(Q)&=\{(I_0,u+4,v),(I_1,p_1,q_1),\ldots,(I_h,p_h,q_h)\},\\
 P_Y(Q)&=\{(I_0,u,v+4),(I_1,p_1,q_1),\ldots,(I_h,p_h,q_h)\},\\
 P_j(Q)&=\{(I_0\sqcup I_j,u+p_j,v+q_j)\}
          \cup\{(I_i,p_i,q_i):i\ne0,j\}.
\end{aligned}
 \label{eq:des-pointing-faces}
\end{equation}
Put
\begin{equation}
\begin{gathered}
 \beta_X(j)=\frac{p_j}{u+p_j},\qquad
 \beta_Y(j)=\frac{q_j}{v+q_j},\\
 A_Q=\frac{1+u+\sum_jp_j}{u+4},\qquad
 B_Q=\frac{1+v+\sum_jq_j}{v+4}.
\end{gathered}
 \label{eq:des-pointing-degrees}
\end{equation}
A zero numerator gives the value zero in the first line.  For a balanced
decomposition $P$, introduce signed rational cellular coefficients
$(a_P,b_P,c_P,d_P)$, summing to one, at the $X$-flag, $Y$-flag, root flag,
and second distinguished marking of its moving circle.  They are not
multisection weights.

The \emph{one-marking pointing chain} uses only the first distinguished marking: set $d_P=0$
and define, inductively on the number of components,
\begin{equation}
 R_Q^{(1)}=1+\sum_{j=1}^h
 \bigl(\beta_X(j)a_{P_j(Q)}+\beta_Y(j)b_{P_j(Q)}\bigr).
 \label{eq:des-qone-gauge}
\end{equation}
The \emph{two-marking pointing chain} also retains $j_2\ne j_1$: put
$(a_P,b_P,c_P,d_P)=(0,0,0,1)$ when $j_1$ and $j_2$ lie in the same block.
Otherwise set $d_P=0$ and
\begin{equation}
 R_Q^{(2)}=1+\sum_{\substack{1\leq j\leq h\\j_2\notin I_j}}
 \bigl(\beta_X(j)a_{P_j(Q)}+\beta_Y(j)b_{P_j(Q)}\bigr).
 \label{eq:des-qtwo-gauge}
\end{equation}
In either case impose
$A_Qa_{P_X(Q)}+B_Qb_{P_Y(Q)}=R_Q^{(\ell)}$ by the symmetric solution
\begin{equation}
 a_{P_X(Q)}=\frac{A_QR_Q^{(\ell)}}{A_Q^2+B_Q^2},
 \qquad
 b_{P_Y(Q)}=\frac{B_QR_Q^{(\ell)}}{A_Q^2+B_Q^2},
 \qquad c_P=1-a_P-b_P.
 \label{eq:des-pointing-gauge}
\end{equation}
The right sides involve one fewer detached component.  The root block
determines $Q$, so no new variable occurs in two equations at the same level.
Thus \eqref{eq:des-pointing-gauge} defines rational, permutation-natural
coefficients without using a GKT coefficient.  They need not be positive;
for example, $h=u=v=0$ gives $(a,b,c)=(2,2,-3)$.  This is permitted for a
rational cellular chain.

For the one-marking chain, the resulting contraction identity is
\begin{equation}
 \sum_{j=1}^h
 \bigl(\beta_X(j)a_{P_j(Q)}+\beta_Y(j)b_{P_j(Q)}\bigr)
 -A_Qa_{P_X(Q)}-B_Qb_{P_Y(Q)}=-1.
 \label{eq:des-qone-row}
\end{equation}
For the two-marking chain, if $j_2\notin I_0$, the same difference omits the unique block
containing $j_2$ and the marked-pointing branch contributes $+1$; the sum is
$-1+1=0$.  If $j_2\in I_0$, the defining same-component case sets all these
coefficients to zero.

We now realize these identities as actual cellular boundaries.  List the named
points of the oriented moving circle in cyclic order as
$p_0,p_1,\ldots,p_{m-1}$, with $p_0$ the root direction, and let $e_i$ be
the arc from $p_i$ to $p_{i+1}$, indices modulo $m$.  The one-marking signed
zero-chain is
\begin{equation}
 z_Q^{(1)}=p_0
 +\sum_{j=1}^h
  (\beta_X(j)a_{P_j(Q)}+\beta_Y(j)b_{P_j(Q)})p_j
 -A_Qa_{P_X(Q)}p_X-B_Qb_{P_Y(Q)}p_Y,
 \label{eq:des-qone-zero-chain}
\end{equation}
where coincident named points are combined.  Equation
\eqref{eq:des-qone-row} says that its total coefficient is zero.  If
$j_2\notin I_0$, define $z_Q^{(2)}$ by omitting the terms with
$j_2\in I_j$ and adding the marked point with coefficient one; if
$j_2\in I_0$, put $z_Q^{(2)}=0$.  The preceding two-marking calculation again gives
total coefficient zero.

For any zero-chain $z=\sum_i\zeta_ip_i$ of total coefficient zero, set
\begin{equation}
 \mathcal P_{p_0}(z)=\sum_{j=1}^{m-1}
   \left(-\sum_{i=1}^{j}\zeta_i\right)e_j,
 \qquad \partial\mathcal P_{p_0}(z)=z.
 \label{eq:des-moving-circle-primitive}
\end{equation}
The boundary formula follows from $\partial e_i=p_{i+1}-p_i$.  It is the
unique primitive whose coefficient on the root arc is zero.  Compatibility
with the complete face diagram is the following direct calculation:
\begin{equation}
 F_*z_Q^{(\ell)}=z_{FQ}^{(\ell)},\qquad
 F_*\mathcal P_{p_0}(z_Q^{(\ell)})
 =\mathcal P_{F(p_0)}(z_{FQ}^{(\ell)}).
 \label{eq:des-pointing-face-naturality}
\end{equation}
For a collision, coincident named points are combined and both adjacent arc
coefficients are repeated.  For a forest contraction, substitution of
\eqref{eq:des-qone-gauge} or \eqref{eq:des-qtwo-gauge} is exactly the
coefficient of the contracted point.  A closed splitting deletes its named
block and leaves the lower-support formula occurring on the right side of the
same recursion.  Wall and affine-overlap maps preserve the pointing flags and
act only on the transported coefficient.  Stabilization deletes a zero
coefficient at an unstable flag, while forgetting and relabelling merely
reindex the cyclic list.  These are all cases in
Proposition~\ref{prop:des-boundary-stratification}, so they prove
\eqref{eq:des-pointing-face-naturality}; it is not inferred from subdivision
naturality alone.  Summing the primitives over all critical decompositions
gives compatible moving-circle chains
$\mathfrak p^{(1)}$ and $\mathfrak p^{(2)}$.  Put
\begin{equation}
 \Psi^{(\ell)}_{I,\mathbf d,N}
 :=\Psi_{I,\mathbf d,N}[\mathfrak p^{(\ell)}],\qquad \ell=1,2.
 \label{eq:des-two-pointing-chains}
\end{equation}
For each named flag $w$, the corresponding cell is fibered over the actual
cotangent zero cobordism $\mathcal Z_{\delta,w}$ from
\eqref{eq:des-cotangent-zero-boundary}.  The coefficients sum to one, so the
$u_\psi=1$ facets combine to one copy of the
degree-$\mathbf d+\mathbf e_{j_1}$ Euler chain.  The positive-weight PL
perturbations are the same in both constructions; only the signed rational
coefficients on the separate cotangent cobordisms differ.

On a contraction facet, the exact-data bundle map restricts $E_\delta^-$ to
the lower-descendent Euler chain $\mathcal E_{I,\mathbf d}$, while the degree
of the pointing map is the corresponding coefficient $\beta_X(j)$,
$\beta_Y(j)$, $A_Q$, or $B_Q$ in
\eqref{eq:des-pointing-degrees}.  Applying the boundary pushforward and then
$\mu_\Gamma$ therefore multiplies every term in
\eqref{eq:des-qone-row} by the common scalar
$\mathcal A_{\mathrm{ann}}(I,\mathbf d)$.  The resulting contraction moment
is $-\mathcal A_{\mathrm{ann}}(I,\mathbf d)$ for the one-marking chain and
zero for the two-marking chain.  Outward-normal-first orientation reverses
this induced moving-circle boundary.  Hence the geometric $\partial_{\psi}$ term in
\eqref{eq:des-absolute-boundary} has one- and two-marking moments
$+\mathcal A_{\mathrm{ann}}(I,\mathbf d)$ and zero, respectively.

Finally, the normalized radial thimble restricts to a zero-dimensional
closed factor exactly when the Euler rank
$\ell_A+m_A+\sum_{i\in A}d_i$ equals the closed dimension $|A|-2$,
which is equivalent to $N_A=-1$.  Equations
\eqref{eq:des-thimble} and \eqref{eq:des-closed-factor} then give coefficient
$4\cdot\frac14=1$ and the product
$\mathcal C_A^{\mathrm{ext}}\mathcal A_{\mathrm{ann}}(Q_A,
\mathbf d|_{I\setminus A})$.  If the ranks do not agree, the closed Euler
pushforward has nonzero residual degree; coefficient-order terms lie in
$\mathfrak m_S^N$.  These are precisely the residual-degree faces, and their
degree-zero moment is zero.

\subsubsection{Stokes evaluation and support reduction}

Apply the contact cocycle to the two explicitly constructed chains
\eqref{eq:des-moving-circle-primitive}.  The degree-zero values of every
term in \eqref{eq:des-absolute-boundary} are
\begin{center}
\begin{tabularx}{\textwidth}{@{}
 >{\raggedright\arraybackslash}p{.14\textwidth}
 >{\raggedright\arraybackslash}p{.40\textwidth}
 >{\raggedright\arraybackslash}X@{}}
\toprule
Face & one-marking value & two-marking value\\
\midrule
$\partial_{\mathrm{can}}$ &
$\mathcal A_{\mathrm{ann}}(I,\mathbf d+\mathbf e_{j_1})$ & same\\
$\partial_{\mathrm{oc}}$ & threshold sum over $j_1\in A$ & threshold sum over $j_1\in A$, $j_2\notin A$\\
$\partial_{\mathrm{wall}}$ & $0$ by \eqref{eq:des-wc-boundary} & $0$\\
$\partial_{\psi}$ & $\mathcal A_{\mathrm{ann}}(I,\mathbf d)$ & $0$\\
$\partial_{\mathrm{exch}}$ & $0$ by orientation reversal & $0$\\
$\partial_{\mathrm{deg}}$ & $0$ by degree or truncation & $0$\\
\bottomrule
\end{tabularx}
\end{center}
Since the contact trace kills a total boundary, this table and the sign
$\partial_{\mathrm{can}}-\partial_{\mathrm{oc}}$ give
\begin{equation}
\begin{aligned}
 \mathcal A_{\mathrm{ann}}(I,\mathbf d+\mathbf e_{j_1})
 ={}&\sum_{\substack{A\subseteq I,\ j_1\in A,\ |A|\geq2\\N_A=-1}}
 \mathcal C_A^{\mathrm{ext}}
 \mathcal A_{\mathrm{ann}}(Q_A,\mathbf d|_{I\setminus A})
 -\mathcal A_{\mathrm{ann}}(I,\mathbf d),
\end{aligned}
\label{eq:des-qone}
\end{equation}
and, for $j_2\ne j_1$,
\begin{equation}
\begin{aligned}
 \mathcal A_{\mathrm{ann}}(I,\mathbf d+\mathbf e_{j_1})
 ={}&\sum_{\substack{A\subseteq I,\ j_1\in A,\ j_2\notin A,
                         \ |A|\geq2\\N_A=-1}}
 \mathcal C_A^{\mathrm{ext}}
 \mathcal A_{\mathrm{ann}}(Q_A,\mathbf d|_{I\setminus A}).
\end{aligned}
\label{eq:des-qtwo}
\end{equation}
The vacuum value is the empty broken-line chain.  For a singleton at
descendent degree zero, the normalized root thimble followed by
\eqref{eq:des-contact-traces} gives one.  Since no subset of a singleton has
cardinality at least two, \eqref{eq:des-qone} then gives
$\mathcal A_{\mathrm{ann}}(\{i\},d_i+1)
=-\mathcal A_{\mathrm{ann}}(\{i\},d_i)$.  Hence the initial values are
\begin{equation}
 \mathcal A_{\mathrm{ann}}(\varnothing,0)=1,
 \qquad \mathcal A_{\mathrm{ann}}(\{i\},d_i)=(-1)^{d_i}.
 \label{eq:des-initial-values}
\end{equation}
Subtracting \eqref{eq:des-qtwo} from \eqref{eq:des-qone} yields
\begin{equation}
 \mathcal A_{\mathrm{ann}}(I,\mathbf d)
 =\sum_{\substack{A\subseteq I,\ j_1,j_2\in A\\N_A=-1}}
 \mathcal C_A^{\mathrm{ext}}
 \mathcal A_{\mathrm{ann}}(Q_A,\mathbf d|_{I\setminus A}).
 \label{eq:des-support-reduction}
\end{equation}
Every open factor on the right has support
$|Q_A|=|I|-|A|+1\leq|I|-1$.  Hence the recursions and
\eqref{eq:des-initial-values} determine the moment family uniquely.  In
particular, no multi-marking primary open invariant is initial data.

\subsection{GKT comparison and the inverse system}

\subsubsection{Moment comparison and endpoint normal forms}

The two GKT cotangent recursions
\cite[Thm.~5.10]{grossetal2022open} are
\eqref{eq:des-qone}--\eqref{eq:des-qtwo} for $r=s=4$, and their singleton
condition is \eqref{eq:des-initial-values}
\cite[Cor.~5.11(1)]{grossetal2022open}.  Their closed factors are the same
compact closed $W$-spin integrals as \eqref{eq:des-closed-factor}.  Uniqueness
therefore gives
\begin{equation}
 \mathcal A_{\mathrm{ann}}(I,\mathbf d)
 =\mathcal A_{\mathrm{GKT}}(I,\mathbf d).
 \label{eq:des-moment-comparison}
\end{equation}

For a nonempty support $J$, define $r_J,s_J,\ell_J,m_J$ by
\[
 \sum_{i\in J}a_i=r_J+4\ell_J,\qquad
 \sum_{i\in J}b_i=s_J+4m_J,\qquad 0\leq r_J,s_J\leq3,
\]
and put
\begin{equation}
 N_J=\ell_J+m_J-|J|+1+\sum_{i\in J}d_i.
 \label{eq:des-balanced-length}
\end{equation}
The support is active when $N_J\geq0$.  GKT Notation~2.40 then lists the
balanced graphs
$\Gamma_{J,0},\ldots,\Gamma_{J,N_J}$ with profiles
\begin{equation}
 (k_1,k_2)(\Gamma_{J,p})
 =(r_J+4p,s_J+4(N_J-p)).
 \label{eq:des-balanced-profiles}
\end{equation}
Write their moment in the triangular form
\begin{equation}
 \mathcal A_{\mathrm{GKT}}(J,\mathbf d;\nu)
 =\sum_{p=0}^{N_J}q_{J,p}\nu_{\Gamma_{J,p},\mathbf d}
  +P_{J,\mathbf d}(\nu|_{\mathrm{proper\ supports}}),
 \label{eq:des-gkt-moment}
\end{equation}
where
\begin{equation}
 q_{J,p}=
 \frac{\Gamma((r_J+4p+1)/4)\Gamma((s_J+4(N_J-p)+1)/4)}
 {\Gamma((r_J+1)/4)\Gamma((s_J+1)/4)}\ne0.
 \label{eq:des-gkt-diagonal}
\end{equation}
Here $P_{J,\mathbf d}$ is the sum over partitions with at least two blocks,
and therefore involves only proper supports.  Put
\[
 b_{J,\mathbf d}=
 \begin{cases}(-1)^{d_i},&J=\{i\},\\0,&|J|\geq2.\end{cases}
\]
Starting with the empty and singleton values, define two labelled systems
recursively on $|J|$ by
\begin{equation}
\begin{aligned}
 &\nu^-_{\Gamma_{J,p},\mathbf d}=0\quad(p>0),&
 q_{J,0}\nu^-_{\Gamma_{J,0},\mathbf d}
 &=b_{J,\mathbf d}-P_{J,\mathbf d}(\nu^-),\\
 &\nu^+_{\Gamma_{J,p},\mathbf d}=0\quad(p<N_J),&
 q_{J,N_J}\nu^+_{\Gamma_{J,N_J},\mathbf d}
 &=b_{J,\mathbf d}-P_{J,\mathbf d}(\nu^+).
\end{aligned}
 \label{eq:des-extreme-systems}
\end{equation}
The Gamma factors in \eqref{eq:des-gkt-diagonal} are nonzero, and every term
in $P_{J,\mathbf d}$ has smaller support.  Induction therefore proves
existence and uniqueness.  Substitution in \eqref{eq:des-gkt-moment} gives
the required singleton moment $(-1)^{d_i}$ and zero higher-support moment;
the construction is twist-preserving permutation invariant and has the
prescribed vanishing outside its single extreme graph.  These are precisely
the chamber-index conditions of \cite[Def.~3.47]{grossetal2022open}.
GKT Theorem~5.13 and Corollary~5.14 consequently realize the two systems by
symmetric canonical multisections.

It remains to show that the annular endpoints give these particular systems.
For a fixed active support $J$, the balanced exponent line has the integral
coordinate
\begin{equation}
 \vartheta_J(r_J+4p,s_J+4(N_J-p))=p.
 \label{eq:des-slope-coordinate}
\end{equation}
The critical wall between consecutive profiles $p-1$ and $p$ is the
hyperplane $\vartheta_J=p-\tfrac12$; this is
\eqref{eq:des-critical-dictionary} below written in the affine cofactor
coordinate.  The stops defining the corridor lie in the two unbounded
chambers $\vartheta_J<\tfrac12$ and
$\vartheta_J>N_J-\tfrac12$.  Substituting these inequalities in
\eqref{eq:des-position-polytope} shows that a one-block cell meeting the
incoming facet has $p=0$, while one meeting the outgoing facet has
$p=N_J$.  Thus the chain maps \eqref{eq:des-endpoint-maps} kill all other
one-block profiles before the contact moment is taken.

For a fixed active $J$, evaluation of the retained connected cell has
diagonal weight $q_{J,0}$ or $q_{J,N_J}$ by
\eqref{eq:des-contact-traces}; every cut-forest cell gives exactly
$P_{J,\mathbf d}$ by \eqref{eq:des-factorization-cocycle}.  The scalar
comparison \eqref{eq:des-moment-comparison} therefore turns the two endpoint
equations into \eqref{eq:des-extreme-systems}.  Induction on $|J|$ identifies
the endpoint coefficients with $\nu^-$ and $\nu^+$.  Since GKT moments are
constant across chamber indices
\cite[Cor.~4.28]{grossetal2022open}, this argument cannot and does not
identify an unspecified chamber representative.

For reference, the shifted exponent of GKT's critical graph
$\Lambda_{J,p}$ is
\begin{equation}
 (k_1-1,k_2-1)(\Lambda_{J,p})
 =(r_J+4(p-1),s_J+4(N_J-p)),
 \label{eq:des-critical-dictionary}
\end{equation}
which places it precisely on the separating hyperplane used above.  Its
vector field is the one whose action on $W_0$ is
\eqref{eq:des-wc-boundary}.  No equality of a raw annular wall coefficient
with a GKT chamber coefficient is used or asserted.

\subsubsection{The inverse-system class and completion of the proof}

We finally define the chain asserted in
Theorem~\ref{thm:intro-descendent}.  For $\delta\in\mathcal E$, let
$I(\delta)$ and $\mathbf d(\delta)$ be its labelled support and descendent
vector, and set
\begin{equation}
 \mathsf C^{\mathrm{lab}}_{S,N}
 =\bigoplus_{\delta\in\mathcal E}
   \mathsf C^{\mathrm{tot}}_{I(\delta),\mathbf d(\delta),N}.
 \label{eq:des-window-total-complex}
\end{equation}
This direct sum has the literal restriction maps of the finite window.

Let $\mathcal A_S$ be the finite set of twist-descendent label types
$\alpha=(a,b,d)$ appearing in $S$, and let $\Gamma_{\QQ}(\mathcal A_S)$ be
the divided-power algebra with basis $x_\alpha^{[m]}$ and multiplication
$x_\alpha^{[m]}x_\alpha^{[n]}
=\binom{m+n}{m}x_\alpha^{[m+n]}$.  For a multiplicity vector
$\mathbf m=(m_\alpha)$, put
$G_{\mathbf m}=\prod_\alpha\mathfrak S_{m_\alpha}$.  Let
$\mathcal E[\mathbf m]\subset\mathcal E$ be the exact data whose labelled
supports contain exactly $m_\alpha$ labels of each type, and define
\begin{equation}
 \mathsf C^{\mathrm{lab}}_{S,N}[\mathbf m]
 =\bigoplus_{\delta\in\mathcal E[\mathbf m]}
   \mathsf C^{\mathrm{tot}}_{I(\delta),\mathbf d(\delta),N}.
 \label{eq:des-labelled-multiplicity-summand}
\end{equation}
The group $G_{\mathbf m}$ acts by permuting equal-type leaf labels on the
forest, open and closed exact data, moduli factors, determinant lines,
pointing sections, and coefficient monomials.  The window contains every
such relabelling.  Each boundary identification and every map
$F_{\Lambda\Gamma}$ and $\phi_{\Lambda\Gamma}$ is equivariant; therefore
the action is by chain automorphisms of
\eqref{eq:des-labelled-multiplicity-summand}.  Define the
coordinate-symmetric target complex
\begin{equation}
 \mathsf C^{\mathrm{dp}}_{S,N}
 =\bigoplus_{\mathbf m}
   (\mathsf C^{\mathrm{lab}}_{S,N}[\mathbf m])_{G_{\mathbf m}}
   \otimes\prod_\alpha x_\alpha^{[m_\alpha]}.
 \label{eq:des-divided-power}
\end{equation}
On an oriented labelled cell generator $c$ with stabilizer
$\operatorname{Stab}(c)\subset G_{\mathbf m}$, define
\begin{equation}
 p_{\mathrm{dp}}(c)
 =\sum_{g\in G_{\mathbf m}/\operatorname{Stab}(c)}[gc]
   \otimes\prod_\alpha x_\alpha^{[m_\alpha]},
 \label{eq:des-dp-pushforward}
\end{equation}
and extend linearly after choosing one generator from each orbit.  This is a
chain map because the differential is equivariant and preserves
$\mathbf m$.  Under the power-series realization
$x_\alpha^{[m]}\mapsto x_\alpha^m/m!$, the scalar multiplying the unlabelled
orbit is
\begin{equation}
 \frac{|G_{\mathbf m}|/|\operatorname{Stab}(c)|}
      {\prod_\alpha m_\alpha!}
 =\frac1{|\operatorname{Stab}(c)|},
 \label{eq:des-orbit-stabilizer-factor}
\end{equation}
which is the required reciprocal automorphism factor.  The same computation
applies to every degeneration face because its stabilizer is computed after
the corresponding labelled boundary map.

Define
\begin{equation}
 \BL_{S,N}
 =p_{\mathrm{dp}}
   \bigl((\Psi^{(1)}_{I(\delta),\mathbf d(\delta),N})_{
          \delta\in\mathcal E}\bigr)
 \in\mathsf C^{\mathrm{dp}}_{S,N}.
 \label{eq:des-broken-line-class}
\end{equation}
Its relative bordism class is the element
$[\BL_{S,N}]_{\mathrm{PL}}\in\Omega^{\mathrm{PL}}_{S,N}$ of
Theorem~\ref{thm:intro-descendent}.
Root transport and factorization decorations are retained in
\eqref{eq:des-broken-line-class}; the contact moment is applied only
afterward.

The naturality in Theorem~\ref{thm:des-absolute-boundary} makes
\eqref{eq:des-window-total-complex}--\eqref{eq:des-broken-line-class} an
inverse system over windows, deformation orders, and finite root stages.
This is not an extension inferred from the scalar GKT comparison: the
multisections, wall cylinders, corolla primitives, and root thimbles on a
smaller window are literal restrictions of those on a larger window.
Every monomial coefficient of the formal genus-zero descendent potential
belongs to some finite divisor- and dependency-closed window, and framed
transport and $K$ act before scalarization.  The inverse system therefore
is the $K$-equivariant inverse system of theta-decorated relative PL bordism
classes asserted in Theorem~\ref{thm:intro-descendent}; applying the contact
moment coefficientwise gives the descendent potential.  The construction
does not identify this bordism module with the theta algebra.

\begin{proof}[Completion of the proof of Theorem~\ref{thm:intro-descendent}]
The boundary theorem, Theorem~\ref{thm:des-absolute-boundary}, constructs the
relative PL chain and identifies all of its exterior faces.  Applying the
contact cocycle gives \eqref{eq:des-qone}--\eqref{eq:des-qtwo}; together with
\eqref{eq:des-initial-values}, these relations determine the moments by
induction on the support.  Equation~\eqref{eq:des-moment-comparison} identifies
them with the GKT moments, while \eqref{eq:des-extreme-systems} identifies the
two annular restrictions.  Finally, \eqref{eq:des-broken-line-class} and the
literal restriction maps just established form the asserted inverse system.
These statements give every construction, recursion, endpoint identification,
and naturality assertion in Theorem~\ref{thm:intro-descendent}.
\end{proof}
